\chapter{Chapter06 — Mod Three}
%%% generated by the blueprint pipeline from FABL.Chapter06.Pseudorandomness.ModThree
\section{Overview}

Declarations of the Lean module \texttt{FABL.Chapter06.Pseudorandomness.ModThree}.

\section{Declarations}

\begin{definition}[Integer sum of sign coordinates]
\lean{FABL.signCoordinateIntegerSum}
\label{FABL.signCoordinateIntegerSum}
\leanok
For a sign vector $x$ of length $n$, its integer coordinate sum is $\sum_{i=1}^n x_i$,
with each coordinate regarded as an integer.
\end{definition}

\begin{definition}[Divisibility Boolean function modulo three]
\lean{FABL.modThreeBoolean}
\statementsource{mit-18405-razborov-smolensky}{pdf-d3f464808105-p001-b004}
\statementsource{goldreich-approximation-method}{pdf-988dc9f0ae91-p004-b004}
\label{FABL.modThreeBoolean}
\leanok
\uses{FABL.ZeroOneFunction, FABL.signCoordinateIntegerSum}
For a nonnegative integer $n$, define a Boolean function on the $n$-dimensional sign
cube that returns true at $x$ exactly when $3$ divides the integer sum of the sign
coordinates of $x$.
\end{definition}

\begin{definition}[Real-valued mod-three quantity]
\lean{FABL.modThree}
\statementsource{mit-18405-razborov-smolensky}{
  pdf-d3f464808105-p001-b004,
  pdf-d3f464808105-p002-b003
}
\statementsource{goldreich-approximation-method}{pdf-988dc9f0ae91-p011-b002}
\label{FABL.modThree}
\leanok
\uses{FABL.ZeroOneFunction.toReal, FABL.modThreeBoolean}
For each natural number $n$, define a function on sign vectors of length $n$ by
assigning to each vector the real value of its associated mod-three Boolean quantity.
\end{definition}

\begin{theorem}[Mod three indicator formula]
\lean{FABL.modThree_apply}
\label{FABL.modThree_apply}
\leanok
\uses{FABL.ZeroOneFunction.toReal, FABL.modThree, FABL.modThreeBoolean, FABL.signCoordinateIntegerSum, FABL.zeroOneValue}
\difficulty{1}
Let $n$ be a nonnegative integer and let $x\in\{-1,1\}^n$. The function that is $1$
when the integer sum of the coordinates of $x$ is divisible by $3$, and $0$ otherwise,
has value $1$ at $x$ if that divisibility holds and value $0$ otherwise.
\end{theorem}

\begin{definition}[Mod-three coordinate sum]
\lean{FABL.modThreeResidue}
\statementsource{goldreich-approximation-method}{pdf-988dc9f0ae91-p011-b002}
\label{FABL.modThreeResidue}
\leanok
For a natural number $n$ and a sign vector $x$ of length $n$, define its residue modulo
$3$ to be the sum of its coordinates, each regarded as an integer and then reduced
modulo $3$.
\end{definition}

\begin{theorem}[Residue divisibility criterion]
\lean{FABL.modThreeResidue_eq_zero_iff}
\label{FABL.modThreeResidue_eq_zero_iff}
\leanok
\uses{FABL.modThreeResidue, FABL.signCoordinateIntegerSum}
\difficulty{3}
Let $n$ be a nonnegative integer and let $x\in\{-1,1\}^n$. The sum of the coordinates
of $x$, reduced modulo $3$, is zero if and only if the integer sum of the coordinates is
divisible by $3$.
\end{theorem}

\begin{theorem}[Residue indicator formula]
\lean{FABL.modThree_eq_residueIndicator}
\label{FABL.modThree_eq_residueIndicator}
\leanok
\uses{FABL.modThree, FABL.modThreeResidue, FABL.modThreeResidue_eq_zero_iff, FABL.modThree_apply, FABL.signCoordinateIntegerSum}
\difficulty{4}
Let $n$ be a nonnegative integer and let $x\in\{-1,1\}^n$. The indicator that the
integer sum of the coordinates of $x$ is divisible by $3$ equals $1$ when the sum
reduced modulo $3$ is zero, and equals $0$ otherwise.
\end{theorem}

\begin{definition}[Primitive cubic root of unity]
\lean{FABL.modThreeCubicRoot}
\label{FABL.modThreeCubicRoot}
\leanok
Let $\omega=\exp(2\pi i/3)$. This defines a complex primitive cube root of unity.
\end{definition}

\begin{theorem}[Cartesian form of the cubic root]
\lean{FABL.modThreeCubicRoot_eq}
\label{FABL.modThreeCubicRoot_eq}
\leanok
\uses{FABL.modThreeCubicRoot}
\difficulty{4}
For $\zeta=\exp(2\pi i/3)$, one has $\zeta=-\frac12+\frac{\sqrt3}{2}i$, where
$i=\sqrt{-1}$.
\end{theorem}

\begin{theorem}[Square of the cubic root]
\lean{FABL.modThreeCubicRoot_sq}
\label{FABL.modThreeCubicRoot_sq}
\leanok
\uses{FABL.modThreeCubicRoot}
\difficulty{4}
For $\zeta=\exp(2\pi i/3)$, one has $\zeta^2=-\frac12-\frac{\sqrt3}{2}i$, where
$i=\sqrt{-1}$.
\end{theorem}

\begin{theorem}[Cubic root identity]
\lean{FABL.modThreeCubicRoot_cube}
\label{FABL.modThreeCubicRoot_cube}
\leanok
\uses{FABL.modThreeCubicRoot}
\difficulty{3}
For $\zeta=\exp(2\pi i/3)$, one has $\zeta^3=1$.
\end{theorem}

\begin{definition}[Additive character modulo three]
\lean{FABL.modThreeCubicCharacter}
\label{FABL.modThreeCubicCharacter}
\leanok
\uses{FABL.modThreeCubicRoot_cube}
Let $\omega=\exp(2\pi i/3)$. The associated additive character on $\bbZ/3\bbZ$ is the
character that sends the residue class of $1$ to $\omega$.
\end{definition}

\begin{theorem}[Character value at zero]
\lean{FABL.modThreeCubicCharacter_zero}
\label{FABL.modThreeCubicCharacter_zero}
\leanok
\uses{FABL.modThreeCubicCharacter}
\difficulty{1}
Let $\zeta=\exp(2\pi i/3)$, and let $\chi$ be the additive character of $\mathbb
Z/3\mathbb Z$ determined by $\chi(1)=\zeta$. Then $\chi(0)=1$.
\end{theorem}

\begin{theorem}[Cubic character at one]
\lean{FABL.modThreeCubicCharacter_one}
\label{FABL.modThreeCubicCharacter_one}
\leanok
\uses{FABL.modThreeCubicCharacter, FABL.modThreeCubicRoot}
\difficulty{3}
Let $\chi$ be the character on $\mathbb Z/3\mathbb Z$ that sends $1$ to
$\zeta=\exp(2\pi\mathrm{i}/3)$. Then $\chi(1)=\zeta$.
\end{theorem}

\begin{theorem}[Character value at residue two]
\lean{FABL.modThreeCubicCharacter_two}
\label{FABL.modThreeCubicCharacter_two}
\leanok
\uses{FABL.modThreeCubicCharacter, FABL.modThreeCubicRoot}
\difficulty{3}
Let $\zeta=\exp(2\pi i/3)$, and let $\chi$ be the additive character of $\mathbb
Z/3\mathbb Z$ determined by $\chi(1)=\zeta$. Then $\chi(2)=\zeta^2$.
\end{theorem}

\begin{theorem}[Cubic character at minus one]
\lean{FABL.modThreeCubicCharacter_neg_one}
\label{FABL.modThreeCubicCharacter_neg_one}
\leanok
\uses{FABL.modThreeCubicCharacter, FABL.modThreeCubicRoot}
\difficulty{3}
Let $\chi$ be the character on $\mathbb Z/3\mathbb Z$ that sends $1$ to
$\zeta=\exp(2\pi\mathrm{i}/3)$. Then $\chi(-1)=\zeta^2$.
\end{theorem}

\begin{theorem}[Cubic character filter identity]
\lean{FABL.modThreeCubicFilter}
\label{FABL.modThreeCubicFilter}
\leanok
\uses{FABL.modThreeCubicCharacter, FABL.modThreeCubicCharacter_one, FABL.modThreeCubicCharacter_two, FABL.modThreeCubicCharacter_zero, FABL.modThreeCubicRoot_eq, FABL.modThreeCubicRoot_sq}
\difficulty{4}
Let $\zeta=\exp(2\pi i/3)$, and let $\chi$ be the additive character of $\mathbb
Z/3\mathbb Z$ determined by $\chi(1)=\zeta$. For every $a\in\mathbb Z/3\mathbb Z$, the
indicator of $a=0$ satisfies $\mathbf
1_{\{0\}}(a)=\frac13+\frac23\operatorname{Re}(\chi(a))$.
\end{theorem}

\begin{definition}[Complex factor from a sign]
\lean{FABL.modThreeCubicFactor}
\label{FABL.modThreeCubicFactor}
\leanok
\uses{FABL.Sign, FABL.signValue}
For a sign $s$, define the complex number $-\frac12 + \frac{\sqrt{3}}2 s i$.
\end{definition}

\begin{theorem}[Cubic character on a sign residue]
\lean{FABL.modThreeCubicCharacter_intCast_sign}
\label{FABL.modThreeCubicCharacter_intCast_sign}
\leanok
\uses{FABL.Sign, FABL.modThreeCubicCharacter, FABL.modThreeCubicCharacter_neg_one, FABL.modThreeCubicCharacter_one, FABL.modThreeCubicFactor, FABL.modThreeCubicRoot_eq, FABL.modThreeCubicRoot_sq, FABL.signValue_neg_one, FABL.signValue_one}
\difficulty{4}
Let $s\in\mathbb Z^\times=\{-1,1\}$, and let $\chi$ be the character on $\mathbb
Z/3\mathbb Z$ that sends $1$ to $\zeta=\exp(2\pi\mathrm{i}/3)$. Then $$\chi(s\bmod
3)=-\frac12+\frac{\sqrt3}{2}s\,\mathrm{i}.$$
\end{theorem}

\begin{theorem}[Additive character of a finite sum]
\lean{FABL.addChar_map_finset_sum_eq_prod}
\label{FABL.addChar_map_finset_sum_eq_prod}
\leanok
\difficulty{4}
Let $A$ be an additive commutative monoid, let $M$ be a commutative monoid, and let
$\chi:A\to M$ be an additive character. For every finite set $s$ and every family
$(a_i)_{i\in\iota}$ in $A$, $\chi\bigl(\sum_{i\in s}a_i\bigr)=\prod_{i\in s}\chi(a_i)$.
\end{theorem}

\begin{theorem}[Cubic character of a sign sum]
\lean{FABL.modThreeCubicCharacter_residue}
\label{FABL.modThreeCubicCharacter_residue}
\leanok
\uses{FABL.addChar_map_finset_sum_eq_prod, FABL.modThreeCubicCharacter, FABL.modThreeCubicCharacter_intCast_sign, FABL.modThreeCubicFactor, FABL.modThreeResidue}
\difficulty{4}
Let $x\in\{-1,1\}^n$, and let $a$ be the sum of its coordinates in $\mathbb Z/3\mathbb
Z$. With $\chi$ the character on $\mathbb Z/3\mathbb Z$ sending $1$ to
$\zeta=\exp(2\pi\mathrm{i}/3)$, one has
$$\chi(a)=\prod_{i=1}^n\left(-\frac12+\frac{\sqrt3}{2}x_i\,\mathrm{i}\right).$$
\end{theorem}

\begin{theorem}[Cubic root filter formula]
\lean{FABL.modThree_eq_cubicRootFilter}
\label{FABL.modThree_eq_cubicRootFilter}
\leanok
\uses{FABL.modThree, FABL.modThreeCubicCharacter_residue, FABL.modThreeCubicFactor, FABL.modThreeCubicFilter, FABL.modThree_eq_residueIndicator}
\difficulty{1}
Let $n$ be a nonnegative integer and let $x\in\{-1,1\}^n$. The indicator that the
integer sum of the coordinates of $x$ is divisible by $3$ equals
$\frac13+\frac23\operatorname{Re}\left(\prod_{i=1}^n\left(-\frac12+\frac{\sqrt3}{2}x_i\,\mathrm{i}\right)\right)$,
where $\mathrm{i}$ is the imaginary unit.
\end{theorem}

\begin{definition}[Complex term from a coordinate set]
\lean{FABL.modThreeComplexTerm}
\label{FABL.modThreeComplexTerm}
\leanok
\uses{FABL.monomial}
For a finite coordinate set $S$ and a sign vector $x$ of length $n$, define
$\bigl(-\sqrt{3}\bigr)^{|S|}\prod_{j\in S}x_j\cdot i^{|S|}$.
\end{definition}

\begin{theorem}[Product of cubic complex factors]
\lean{FABL.modThreeCubicVariableProduct}
\label{FABL.modThreeCubicVariableProduct}
\leanok
\uses{FABL.modThreeComplexTerm, FABL.monomial, FABL.signValue}
\difficulty{4}
Let $n$ be a nonnegative integer, let $S$ be a subset of the coordinate indices
$	ext{Fin}(n)$, and let $x$ be a sign vector in $\{-1,1\}^n$. The product over $i\in
S$ of $(-\sqrt{3}\,x_i)\,\mathrm{i}$ equals $(-\sqrt{3})^{|S|}(\prod_{i\in
S}x_i)\,\mathrm{i}^{|S|}$, where $\mathrm{i}$ is the imaginary unit.
\end{theorem}

\begin{theorem}[Complex product expansion]
\lean{FABL.modThreeCubicProduct_expansion}
\label{FABL.modThreeCubicProduct_expansion}
\leanok
\uses{FABL.modThreeComplexTerm, FABL.modThreeCubicFactor, FABL.modThreeCubicVariableProduct, FABL.signValue}
\difficulty{5}
Let $n$ be a natural number and let $x=(x_0,\ldots,x_{n-1})\in\{−1,1\}^n$. For each
subset $S\subseteq\{0,\ldots,n-1\}$, write $m_S(x)=\prod_{i\in S}x_i$. Then
\[\prod_{i=0}^{n-1}\left(-\frac12+\frac{\sqrt3}{2}x_i\,i\right)=\left(-\frac12\right)^n\sum_{S\subseteq\{0,\ldots,n-1\}}(-\sqrt3)^{|S|}m_S(x)\,i^{|S|},\]
where $i=\sqrt{-1}$.
\end{theorem}

\begin{theorem}[Real part of an even power of the imaginary unit]
\lean{FABL.re_I_pow_eq_modFourPhase}
\label{FABL.re_I_pow_eq_modFourPhase}
\leanok
\difficulty{5}
For every even nonnegative integer $k$, the real part of $i^k$ is $(-1)^{\lfloor (k\bmod
4)/2\rfloor}$.
\end{theorem}

\begin{theorem}[Real part of an odd power of the imaginary unit]
\lean{FABL.re_I_pow_eq_zero_of_odd}
\label{FABL.re_I_pow_eq_zero_of_odd}
\leanok
\difficulty{4}
For every odd nonnegative integer $k$, the real part of $i^k$ is zero.
\end{theorem}

\begin{theorem}[Real part of a mod-three complex term]
\lean{FABL.modThreeComplexTerm_re}
\label{FABL.modThreeComplexTerm_re}
\leanok
\uses{FABL.modThreeComplexTerm, FABL.monomial, FABL.re_I_pow_eq_modFourPhase, FABL.re_I_pow_eq_zero_of_odd}
\difficulty{4}
Let $S\subseteq\{1,\ldots,n\}$ and $x\in\{-1,1\}^n$. The real part of
$(-\sqrt{3})^{|S|}\bigl(\prod_{i\in S}x_i\bigr)\,\mathrm{i}^{|S|}$ is
$$\begin{cases}(-1)^{(|S|\bmod 4)/2}(\sqrt{3})^{|S|}\displaystyle\prod_{i\in
S}x_i,&|S|\text{ is even},\0,&|S|\text{ is odd}.\end{cases}$$
\end{theorem}

\begin{theorem}[Real part of the cubic product]
\lean{FABL.modThreeCubicProduct_re}
\label{FABL.modThreeCubicProduct_re}
\leanok
\uses{FABL.modThreeComplexTerm_re, FABL.modThreeCubicFactor, FABL.modThreeCubicProduct_expansion, FABL.monomial}
\difficulty{3}
Let $n$ be a natural number and let $x=(x_0,\ldots,x_{n-1})\in\{−1,1\}^n$, and write
$m_S(x)=\prod_{i\in S}x_i$ for each subset $S\subseteq\{0,\ldots,n-1\}$. The real part
of the product satisfies
\[\operatorname{Re}\prod_{i=0}^{n-1}\left(-\frac12+\frac{\sqrt3}{2}x_i\,i\right)=\left(-\frac12\right)^n\sum_{\substack{S\subseteq\{0,\ldots,n-1\}\\ |S|\text{ even}}}(-1)^{\lfloor (|S|\bmod 4)/2\rfloor}(\sqrt3)^{|S|}m_S(x),\]
where $i=\sqrt{-1}$.
\end{theorem}

\begin{theorem}[Fourier expansion of the mod three indicator]
\lean{FABL.modThree_fourier_expansion}
\label{FABL.modThree_fourier_expansion}
\leanok
\uses{FABL.modThree, FABL.modThreeCubicProduct_re, FABL.modThree_eq_cubicRootFilter, FABL.monomial}
\difficulty{3}
Let $n$ be a nonnegative integer and let $x\in\{-1,1\}^n$. The indicator that the
integer sum of the coordinates of $x$ is divisible by $3$ equals
$\frac13+\frac23(-\frac12)^n\sum_{S\subseteq\{1,\ldots,n\},\ |S|\text{
even}}(-1)^{(|S|\bmod 4)/2}(\sqrt3)^{|S|}\prod_{i\in S}x_i$.
\end{theorem}

\begin{definition}[Mod-three Fourier coefficient]
\lean{FABL.modThreeFourierCoefficient}
\label{FABL.modThreeFourierCoefficient}
\leanok
For a natural number $n$ and a subset $S\subseteq\{0,\ldots,n-1\}$, define its
coefficient to be $1/3$ when $S$ is empty and $0$ otherwise, plus the following term
when $|S|$ is even (and zero when it is odd):
$\frac{2}{3}(-\frac12)^n(-1)^{\lfloor(|S|\bmod 4)/2\rfloor}(\sqrt{3})^{|S|}$.
\end{definition}

\begin{theorem}[Fourier expansion of the mod-three indicator]
\lean{FABL.modThree_eq_multilinearPolynomial}
\label{FABL.modThree_eq_multilinearPolynomial}
\leanok
\uses{FABL.modThree, FABL.modThreeFourierCoefficient, FABL.modThree_fourier_expansion, FABL.monomial, FABL.multilinearPolynomial}
\difficulty{5}
Let $n$ be a nonnegative integer and let $x\in\{-1,1\}^n$. The function equal to $1$
when the integer sum of the coordinates of $x$ is divisible by $3$, and to $0$
otherwise, is the multilinear polynomial whose coefficient at a subset
$S\subseteq\{1,\ldots,n\}$ is $1/3$ when $S$ is empty, plus
$\frac{2}{3}(-1/2)^n(-1)^{(|S|\bmod 4)/2}(\sqrt{3})^{|S|}$ when $|S|$ is even; this
additional contribution is zero when $|S|$ is odd.
\end{theorem}

\begin{theorem}[Fourier coefficients of the mod-three function]
\lean{FABL.fourierCoeff_modThree}
\label{FABL.fourierCoeff_modThree}
\leanok
\uses{FABL.fourierCoeff, FABL.fourier_expansion_unique, FABL.modThree, FABL.modThreeFourierCoefficient, FABL.modThree_eq_multilinearPolynomial}
\difficulty{3}
Let $n$ be a nonnegative integer and let $S$ be a set of coordinates in the
$n$-dimensional sign cube. The function that is $1$ when the sum of the input signs is
divisible by $3$ and $0$ otherwise has Fourier coefficient at $S$ equal to
$\mathbf{1}_{S=\varnothing}/3$ plus, when $|S|$ is even,
$\frac{2}{3}(-\frac12)^n(-1)^{(|S|\bmod 4)/2}(\sqrt{3})^{|S|}$, and otherwise equal to
just $\mathbf{1}_{S=\varnothing}/3$.
\end{theorem}

\begin{theorem}[Nonempty mod-three coefficients]
\lean{FABL.fourierCoeff_modThree_of_nonempty}
\label{FABL.fourierCoeff_modThree_of_nonempty}
\leanok
\uses{FABL.fourierCoeff, FABL.fourierCoeff_modThree, FABL.modThree, FABL.modThreeFourierCoefficient}
\difficulty{3}
Let $n$ be a nonnegative integer and let $S$ be a nonempty set of coordinates in the
$n$-dimensional sign cube. The Fourier coefficient at $S$ of the indicator that the sum
of the input signs is divisible by $3$ equals $\frac{2}{3}(-\frac12)^n(-1)^{(|S|\bmod
4)/2}(\sqrt{3})^{|S|}$ if $|S|$ is even, and equals $0$ if $|S|$ is odd.
\end{theorem}

\begin{theorem}[Fourier coefficient bound for the mod-three function]
\lean{FABL.abs_fourierCoeff_modThree_le}
\label{FABL.abs_fourierCoeff_modThree_le}
\leanok
\uses{FABL.fourierCoeff, FABL.fourierCoeff_modThree_of_nonempty, FABL.modThree}
\difficulty{4}
Let $n$ be a nonnegative integer, and let $S$ be a nonempty subset of the coordinates of
the $n$-dimensional sign cube. The absolute value of the Fourier coefficient of the
mod-three function on $S$ is at most $\frac{2}{3}\left(\frac{\sqrt{3}}{2}\right)^n$.
\end{theorem}

\begin{theorem}[Fourier regularity of the mod three indicator]
\lean{FABL.modThree_isFourierRegular}
\label{FABL.modThree_isFourierRegular}
\leanok
\uses{FABL.IsFourierRegular, FABL.abs_fourierCoeff_modThree_le, FABL.modThree}
\difficulty{1}
Let $n$ be a nonnegative integer. Every nonempty subset $S\subseteq\{1,\ldots,n\}$
has Fourier coefficient for the indicator of coordinate sums divisible by $3$ with
absolute value at most $\frac23(\frac{\sqrt3}{2})^n$.
\end{theorem}
